\section{Limitations}
\label{sec:limits}

% STE descriptive
\paragraph{The early enrichment ramp is not a progress measure.} The CRT enrichment rises
monotonically from early training in grokked runs. It does the same in runs that never
generalise, and it ends higher than grokked runs reach early (Section~\ref{sec:crt-ramp}).
Neither the ramp's presence, its monotonicity nor its permutation $p$ distinguishes a model
that found the circuit from one that memorised its training split. Only the final enrichment
magnitude does. To settle whether any early signal is diagnostic, a statistic must be scored on
a failed-run control before it is proposed. We now apply that procedure, and we did not apply
it when we first measured the ramp.

\paragraph{The engine's dtype-robustness is unexplained.} Section~\ref{sec:impl} shows that
the engine--PyTorch sparsity gap is precision that acts on PyTorch. It leaves open why our
engine reaches the float64 answer at either dtype. We eliminated three candidates under a
pre-registration with both controls: residual type promotion, matmul accumulation order and
softmax formulation. No candidate survives. An operator-by-operator differential trace of a
single training step at both dtypes would settle it.

\paragraph{No algebraic property predicts grokking time.} Grokking time varies by a factor of
$3.8$ across seeds at one modulus. Two claims of ours about how it orders across moduli were
retracted. One was lost to seed variance, and one to a normalisation that manufactured the order
it reported. Nothing in Table~\ref{tab:moduli} predicts it in this data. To establish
that a property does would need many more seeds per modulus than five, at a cost we could
not meet.

\paragraph{$\omega(n)$ is confirmatory, not a discovery.} The additive-sparsity bands were seen
before the pre-registration that tested them, on the same data. A pre-registration cannot
convert that into a prediction. The confounds are dead and $\omega$ does not displace
zero-divisor density, and this data cannot rank the two (Section~\ref{sec:crt-omega}). Its
designed falsifier at $n = 2^{7}$ landed between the two predicted intervals, and neither
predictor was accurate. To rank the two requires moduli chosen to separate them, trained
after the prediction is fixed.

\paragraph{Three claims rest on a single seed.} The torsor coordinate as first sighted, the
vestigial embedding frequency, and the full mechanistic readout of the addition control are
each one seed. They are marked as such wherever they appear. The first is superseded by a
five-seed result, and nothing rests on it. The other two are reported as observations. The four
additional seeds the runs already support would settle each.

\paragraph{Post-grok instability is observed and not explained.} Under weight decay $1.0$,
accuracy excursions below $0.90$ after grokking are common. There are $19$ across $11$ long
runs, all exactly one logging interval wide. Their rate rises with horizon, $0.33\%$ of
samples at $120{,}000$ steps against $0.09\%$ at $40{,}000$. We previously read one such
excursion, which touched a run's final logged sample, as evidence that a network can lose a
learned circuit. It is an artefact of censored data, and the claim is retracted. What causes the
excursions is not addressed here.

A finer logging interval through one excursion would show whether the circuit or the readout
is perturbed.

\paragraph{One test was underpowered.} A search for a specific rare event returned $0$ of
$453$, where the no-effect expectation was $0.68$. The pre-registered probability of zero
events under the effect was also about $0.51$. So the test distinguishes nothing, and its null
is not reported as a result. Only a longer training horizon can expect enough events to decide
it.

\paragraph{Small models, one task family.} Every result comes from a single-layer
transformer, $d_{\text{model}} = 128$, on $a \cdot b \bmod n$ for $n \le 169$. The weight decay
has one value, and so does the training fraction outside the sweep of
Section~\ref{sec:clock-grok}. We do not show that the stratified mechanism is what larger
models or other algebraic tasks do. \citet{zhong2023clock} establish that architectural and
hyperparameter choices can yield different algorithms on a closely related task. The natural
next test is the same measurement on non-commutative finite groups. There,
\citet{chughtai2023toy}'s representation-theoretic account gives an independent prediction of
what the strata should be.

\paragraph{The earliest three sweeps carry no provenance stamp, and their code version is
recovered by argument rather than by record.} We added the provenance stamp after the first
three sweeps ran. So $37$ artifacts carry no stamp at all and $31$ carry the literal string
\texttt{unknown} in place of a commit. Every later sweep, and every run of the from-scratch
engine, carries a real commit. The affected artifacts are behind
Sections~\ref{sec:causal}, \ref{sec:strata} and \ref{sec:crt}, which is most of what this
paper claims. So we state the recovery explicitly rather than leave it implicit.

For the third sweep, the training script has \emph{exactly one} version in the repository's
history. So its code version is determined, and the \texttt{unknown} string records a stamp
failure rather than an ambiguity. For the second sweep, the script has three versions, and we
check their differences mechanically. The model's forward computation and its
non-instrumented return are byte-identical across all three. The only change is the tuple
returned when activations are requested, and the training and evaluation call sites never
request it.

The first sweep, a six-run scout, ran \emph{before} the repository's first commit, so its
history cannot certify what ran. Instead, the source that the compute platform stored for the
version behind those artifacts is byte-identical (equal SHA-256) to the one committed script.
The training path is therefore determined in all three cases.

What remains true is that this is a reconstruction and not a stamp. It rests on a complete
repository history and, for the scout, on the platform's stored record. Every later artifact
rests on a commit recorded at the moment it was written.

The engine's own first run, the addition gate of Section~\ref{sec:calib}, also predates the
repository and was first saved without a stamp. The engine is deterministic and full-batch,
so we re-ran it at a recorded commit. It reproduced every saved weight, logit, activation and
logged value bit for bit. So its artifacts now carry that commit by reproduction rather than
argument.
% STE end
